\subsection{Toe clearance}
\label{sec:res-clearance}

The camera minimum clearance (lowest toe or heel marker over the obstacle minus the obstacle top) is a safety outcome: a low
trailing clearance is a near-trip. With person-scaled Dot features the clearance is not predictable (Table~\ref{tab:safety}:
$r \approx 0$) and low-clearance crossings are only weakly detectable (AUC 0.63--0.73). The reason is that person scaling removes
each person's own margin, which is most of the clearance. Two further questions were therefore asked (Tables~\ref{tab:clearcells}
and \ref{tab:clearabs}): what sets the clearance from one crossing to the next, and how well the raw IMU foot lift predicts it.

\begin{table}[htbp]
\centering
\caption{Clearance from person-scaled Dot features (random forest): regression and detection of low-clearance crossings.}
\label{tab:safety}
\begin{adjustbox}{max width=\textwidth}
{\footnotesize
\begin{tabular}{p{4.6cm}p{4cm}p{6.2cm}}
\toprule
\textbf{task} & \textbf{inputs} & \textbf{result} \\
\midrule
regression of trail clearance (mm) & trail Dot features, self, forest & MAE 85 mm (predicting the mean: 84); r 0.04; within-person r median 0.28 \\
regression of lead clearance (mm) & lead Dot features, self, forest & MAE 49 mm (predicting the mean: 46); r $-$0.06; within-person r median 0.18 \\
regression of trail clearance (mm) & LEAD Dot features (cueing), self, forest & MAE 90 mm (predicting the mean: 84); r $-$0.14; within-person r median 0.08 \\
detect trailing clearance $\leq$ 0 mm (15 events) & trail Dot features, self, forest & AUC 0.73 \\
detect trailing clearance $\leq$ 30 mm (24 events) & trail Dot features, self, forest & AUC 0.70 \\
detect trailing clearance $\leq$ 50 mm (30 events) & trail Dot features, self, forest & AUC 0.63 \\
\bottomrule
\end{tabular}
}
\end{adjustbox}

\end{table}

\paragraph{What sets the clearance.} Within one participant and one obstacle the clearance still varies from crossing to crossing
(SD 37\,mm leading, 71\,mm trailing). About a third of this variation is explained by how high the toe goes (camera toe peak
above the obstacle: 35\,\% and 34\,\%), and very little by where the foot leaves the floor (take-off distance: 4\,\% and 6\,\%).
The IMU foot lift explains 20\,\% (leading) and 12\,\% (trailing): it follows the camera toe peak only partly ($r = 0.44$ and
0.21), because the sensor sits on the top of the foot and the trailing foot crosses heel-up, so the toe is not where the sensor
is. The provisional rigid-body toe estimate does no better.

\begin{table}[htbp]
\centering
\caption{What sets the crossing-to-crossing clearance (within participant $\times$ obstacle cells): SDs, correlations with the
clearance, and the share of the clearance variance explained.}
\label{tab:clearcells}
\begin{adjustbox}{max width=\textwidth}
{\footnotesize
\begin{tabular}{p{9.4cm}rr}
\toprule
\textbf{within participant x obstacle cells} & \textbf{leading limb} & \textbf{trailing limb} \\
\midrule
crossings (cells) & 260 (60) & 259 (60) \\
clearance SD (mm) & 37 & 71 \\
camera toe peak above the obstacle: SD (mm), r with clearance & 50, r = +0.59 & 84, r = +0.58 \\
camera take-off distance: SD (mm), r with clearance & 142, r = $-$0.19 & 89, r = +0.25 \\
IMU foot lift: SD (mm), r with clearance & 42, r = +0.44 & 37, r = +0.35 \\
IMU toe rise (provisional): SD (mm), r with clearance & 42, r = +0.43 & 34, r = +0.37 \\
r of IMU foot lift with the camera toe peak & +0.44 & +0.21 \\
r of IMU foot lift with the take-off distance & +0.05 & +0.01 \\
share of clearance variance explained: camera toe peak & 35\% & 34\% \\
share of clearance variance explained: camera take-off distance & 4\% & 6\% \\
share of clearance variance explained: camera toe peak + take-off & 37\% & 42\% \\
share of clearance variance explained: IMU foot lift & 20\% & 12\% \\
share of clearance variance explained: IMU toe rise (provisional) & 19\% & 13\% \\
share of clearance variance explained: IMU foot lift + take-off (camera) & 24\% & 18\% \\
\bottomrule
\end{tabular}
}
\end{adjustbox}

\end{table}

\paragraph{Absolute clearance for a new person.} Between people and obstacles the raw IMU foot lift (with the obstacle height) does
predict part of the clearance: mean absolute error 39\,mm for the leading and 65\,mm for the trailing limb, against 49 and 88\,mm for
the training mean (Table~\ref{tab:clearabs}). The camera's own toe peak would reach 34 and 55\,mm. Clearance is thus partly
predictable from the IMU; a trip-risk monitor needs a better toe-height estimate (a measured sensor-to-toe geometry).

\begin{table}[htbp]
\centering
\caption{Absolute clearance of a left-out participant (ordinary least squares): mean absolute error and correlation.}
\label{tab:clearabs}
\begin{adjustbox}{max width=\textwidth}
{\footnotesize
\begin{tabular}{lrr}
\toprule
\textbf{inputs} & \textbf{leading limb: MAE, r} & \textbf{trailing limb: MAE, r} \\
\midrule
constant (training mean) & 49 mm, r = $-$0.71 & 88 mm, r = $-$0.62 \\
IMU: foot lift + obstacle height & 39 mm, r = 0.49 & 65 mm, r = 0.56 \\
IMU + take-off distance (camera) & 42 mm, r = 0.40 & 66 mm, r = 0.55 \\
camera toe peak + take-off distance & 34 mm, r = 0.67 & 55 mm, r = 0.70 \\
\bottomrule
\end{tabular}
}
\end{adjustbox}

\end{table}
